%%%%%%%%%%%%
%
% $Autor: Wings $
% $Datum: 2019-03-05 08:03:15Z $
% $Version: 4250 $
% !TeX spellcheck = en_GB/de_DE
% !TeX encoding = utf8
% !TeX TXS-program:bibliography = txs:///biber
%
%%%%%%%%%%%%

\section{Dataset Preparation}

\textbf{Bounding boxes:}Bounding boxes are the most commonly used type of annotation in computer vision. Bounding boxes are rectangular boxes used to define the location of the target object. They can be determined by the $x$ and $y$ axis coordinates in the upper-left corner and the $x$ and $y$ axis coordinates in the lower-right corner of the rectangle. Bounding boxes are generally used in object detection and localization tasks.
\begin{figure}[H]
	\centering
	\GRAPHICS{0.6}{0.6}{Annotation/boundingbox.png}
	\caption{Bounding box showing coordinates x1,y1,x2,y2,width:w, height: h \cite{Kwon:2021}}
\end{figure} 

\subsection{Labelling Format: YOLO}

In \ac{yolo} labeling format, a \FILE{.txt} file with the same name is created for each image file in the same directory. Each \FILE{.txt} file contains the annotations for the corresponding image file, that is object class, object coordinates, height and width. Below is the format in which this data is stored in a text file for each object and for each image.

\bigskip

<object-class> <x> <y> <width> <height> 

\bigskip

For each object, a new line is created if in an images, multiple objects needs to be annotated. 

\subsection{Annotation of the images}

Annotation tool used: The project \FILE{Yolo{\_}Label}  
is downloaded from the Github repository.\cite{Kwon:2021}
\textbf{Note: }Further information about the annnotation tool and how to use it can be found in the file \FILE{Distortion\_Camera\_Calibration\_OpenCV.tex}in the Concepts folder. 
A \FILE{.exe} file is present which can be run. Firstly, the directory of the images was specified.
Supported image formats are \FILE{.png} and \FILE{.jpg}. Secondly, a text file with the name \FILE{label.txt} containing the clases was specified
This tool creates a \FILE{.txt} file containing the class number, the object coordinates, width and height.As we have 2 different objects: cylinder and cuboid, there are 2 class numbers generated: 0 for cylinder and 1 for cuboid. This order can change depending upon the order of classes we give in the \FILE{label.txt} file. 
The below images show the contents of the text file created for sample images for cylinder and cuboid objects: 

\begin{figure}[H]
	\centering
	\GRAPHICS{1}{1}{Annotation/TxtFileSample1.png}
	\caption{Contents of text file created for a sample image of cuboid}
\end{figure}

\begin{figure}[H]
	\centering
	\GRAPHICS{1}{1}{Annotation/TxtFileSample2.png}
	\caption{Contents of text file created for a sample image of cylinder}
\end{figure}

The text files are saved with the same file name as the image which is annotated.

\begin{figure}[H]
	\centering
	\GRAPHICS{0.6}{0.6}{Annotation/DatasetSampleCylinder.png}
	\caption{Sample of dataset for cylindrical object}
\end{figure}

\begin{figure}[H]
	\centering
	\GRAPHICS{0.8}{0.8}{Annotation/DatasetSampleCuboid.png}
	\caption{Sample of dataset for cuboid object}
\end{figure}

From a database of approx 4000 images taken for both the objects in different angles of the camera and the positions of the objects, 672 images were selected for cylinder and 696 images for cuboid depending upon the clarity of the picture and the background.

\subsection{Challenges}

Tedious to use as we have to manually draw bounding boxes around the objects.

